\documentclass[aspectratio=169,12pt]{beamer}

\usepackage[utf8]{inputenc}
\usepackage[spanish,es-tabla]{babel}
\usepackage[sfdefault]{FiraSans}
\usepackage{graphicx}
\usepackage{booktabs}
\usepackage{tabularx}
\usepackage{hyperref}
\usepackage{xcolor}

\definecolor{unamblue}{RGB}{0,51,102}
\definecolor{unamgold}{RGB}{184,151,43}
\definecolor{cempasuchil}{RGB}{232,112,20}
\definecolor{noche}{RGB}{13,20,38}
\definecolor{crema}{RGB}{249,241,219}

\setbeamercolor{normal text}{fg=noche,bg=crema}
\setbeamercolor{structure}{fg=unamblue}
\setbeamercolor{frametitle}{fg=white,bg=unamblue}
\setbeamercolor{title}{fg=white}
\setbeamercolor{subtitle}{fg=crema}
\setbeamercolor{block title}{fg=white,bg=cempasuchil}
\setbeamercolor{block body}{fg=noche,bg=white}
\setbeamercolor{alerted text}{fg=cempasuchil}
\setbeamertemplate{navigation symbols}{}
\setbeamertemplate{itemize item}{\color{cempasuchil}$\blacktriangleright$}
\setbeamertemplate{itemize subitem}{\color{unamgold}$\bullet$}
\setbeamertemplate{footline}{%
  \leavevmode\hbox{%
    \begin{beamercolorbox}[wd=.82\paperwidth,ht=2.7ex,dp=1.1ex,leftskip=1em]{author in head/foot}%
      \color{white}\scriptsize CEMPASÚCHIL · Equipo 8 · Facultad de Ingeniería, UNAM
    \end{beamercolorbox}%
    \begin{beamercolorbox}[wd=.18\paperwidth,ht=2.7ex,dp=1.1ex,center]{date in head/foot}%
      \color{white}\scriptsize\insertframenumber/\inserttotalframenumber
    \end{beamercolorbox}}}
\setbeamercolor{author in head/foot}{bg=noche}
\setbeamercolor{date in head/foot}{bg=cempasuchil}

\newcommand{\img}[2][]{\includegraphics[#1]{#2}}
\newcommand{\captura}[2]{%
  \begin{minipage}{0.48\linewidth}\centering
    \includegraphics[width=\linewidth,height=0.29\textheight,keepaspectratio]{#1}\\[-0.3em]
    {\scriptsize #2}
  \end{minipage}}

\title{CEMPASÚCHIL}
\subtitle{Ofrenda, memoria y tradición}
\author{Equipo 8}
\institute{Facultad de Ingeniería · Universidad Nacional Autónoma de México\\
Computación Gráfica e Interacción Humano-Computadora · Grupo 04}
\date{Semestre 2027-1}

\begin{document}

\setbeamercolor{background canvas}{bg=noche}
\begin{frame}[plain]
  \vspace{1.2cm}
  \begin{center}
    \includegraphics[width=0.58\textwidth]{../doc/images/cempasuchil.png}\\[0.4cm]
    {\color{crema}\Large Ofrenda, memoria y tradición}\\[0.7cm]
    {\color{white}\normalsize Proyecto final · Equipo 8 · Semestre 2027-1}
  \end{center}
\end{frame}
\setbeamercolor{background canvas}{bg=crema}

\begin{frame}{Concepto de la experiencia}
  \begin{columns}[T,onlytextwidth]
    \column{0.48\textwidth}
      \begin{itemize}
        \item Recorrido nocturno por un panteón mexicano.
        \item Sendero de pétalos, faroles y papel picado hasta una ofrenda de tres niveles.
        \item Siete retratos recuerdan a músicos que murieron a los 27 años.
        \item La escena relaciona memoria, fragilidad y permanencia mediante objetos tradicionales.
      \end{itemize}
    \column{0.49\textwidth}
      \includegraphics[width=\linewidth,height=0.66\textheight,keepaspectratio]{../doc/images/final-general.png}
  \end{columns}
\end{frame}

\begin{frame}{Evolución por entregas}
  \centering
  \captura{../doc/images/etapa-prototipo.png}{Prototipo con primitivas}\hfill
  \captura{../doc/images/etapa-modelado-ofrenda.png}{Modelado completo}\\[0.2cm]
  \captura{../doc/images/etapa-texturizado-ofrenda.png}{Texturizado e iluminación}\hfill
  \captura{../doc/images/final-altar.png}{Integración final}
\end{frame}

\begin{frame}[shrink=10]{Escena y recorrido}
  \begin{columns}[T,onlytextwidth]
    \column{0.62\textwidth}
      \includegraphics[width=\linewidth,height=0.60\textheight,keepaspectratio]{../doc/images/final-camino.png}
    \column{0.35\textwidth}
      \small
      \begin{block}{Trayecto}
        La cámara inicia en la entrada, atraviesa el panteón y sigue el camino de cempasúchil hasta la casa y el altar.
      \end{block}
      \begin{itemize}
        \item Navegación en primera persona
        \item Altura de cámara de 1.3 m
        \item Balanceo suave al caminar
        \item Vuelo opcional para inspección
      \end{itemize}
  \end{columns}
\end{frame}

\begin{frame}{Pipeline gráfico}
  \begin{columns}[T,onlytextwidth]
    \column{0.47\textwidth}
      \begin{block}{Producción de contenido}
        Blender 5.2 para modelado, UV, horneado, materiales, armatures y animaciones externas. Los modelos se intercambian mediante glTF 2.0.
      \end{block}
      \begin{block}{Aplicación}
        C++17, OpenGL 3.3 Core, GLSL 330, GLFW, GLAD, GLM, cgltf y stb\_image.
      \end{block}
    \column{0.50\textwidth}
      \includegraphics[width=\linewidth,height=0.64\textheight,keepaspectratio]{../doc/images/etapa-modelado-alambre.png}
  \end{columns}
\end{frame}

\begin{frame}{Modelado y composición}
  \begin{columns}[T,onlytextwidth]
    \column{0.53\textwidth}
      \captura{../doc/images/detalle-alimentos.png}{Alimentos modelados por el equipo}\hfill
      \captura{../doc/images/detalle-recipientes.png}{Cerámica, vidrio y agua}
    \column{0.43\textwidth}
      \begin{itemize}
        \item 81 modelos cargados
        \item 2,134 nodos de escena
        \item 10,549 instancias
        \item Jerarquías padre--hijo para objetos compuestos
        \item Instanciado para flores, pétalos, veladoras y bardas
      \end{itemize}
  \end{columns}
\end{frame}

\begin{frame}{Texturizado e iluminación}
  \centering
  \captura{../doc/images/etapa-texturizado-pbr.png}{Mapas horneados y materiales}\hfill
  \captura{../doc/images/etapa-fresnel.png}{Fresnel para vidrio y agua}\\[0.25cm]
  \begin{itemize}
    \item Phong para materiales opacos y luces puntuales dinámicas.
    \item Fresnel para transmisión y reflejo en superficies translúcidas.
    \item Skybox nocturno, transparencia alfa y emisión para flamas.
    \item Iluminación fría ambiental y luces cálidas que dirigen la mirada al altar.
  \end{itemize}
\end{frame}

\begin{frame}{Navegación e interacción}
  \begin{columns}[T,onlytextwidth]
    \column{0.48\textwidth}
      \begin{itemize}
        \item \textbf{WASD}: caminar o volar
        \item \textbf{TAB}: capturar o liberar el ratón
        \item \textbf{F}: activar o desactivar vuelo
        \item \textbf{Q/E}: bajar o subir durante el vuelo
        \item \textbf{R}: regresar a la entrada
        \item \textbf{F2}: mostrar mallas
        \item \textbf{L}: encender o apagar faroles
      \end{itemize}
    \column{0.49\textwidth}
      \includegraphics[width=\linewidth,height=0.64\textheight,keepaspectratio]{../doc/images/vista-entrada.png}
  \end{columns}
\end{frame}

\begin{frame}{Animaciones programadas en OpenGL}
  \centering
  \captura{../doc/images/animacion-mariposas.png}{Mariposas: Catmull--Rom y jerarquía cuerpo--alas}\hfill
  \captura{../doc/images/animacion-recorrido-3.png}{Recorrido automático de cámara}\\[0.25cm]
  \captura{../doc/images/animacion-luces-encendidas.png}{Veladoras y luces encendidas}\hfill
  \captura{../doc/images/animacion-luces-apagadas.png}{Veladoras y luces apagadas}
\end{frame}

\begin{frame}{Animaciones creadas en Blender}
  \centering
  \begin{minipage}{0.32\linewidth}\centering
    \includegraphics[width=\linewidth,height=0.44\textheight,keepaspectratio]{../doc/images/animacion-papel-en-movimiento.png}\\
    \small Papel picado con caché de vértices
  \end{minipage}\hfill
  \begin{minipage}{0.32\linewidth}\centering
    \includegraphics[width=\linewidth,height=0.44\textheight,keepaspectratio]{../doc/images/animacion-remolito-petalos-3.png}\\
    \small Remolino de 720 pétalos
  \end{minipage}\hfill
  \begin{minipage}{0.32\linewidth}\centering
    \includegraphics[width=\linewidth,height=0.44\textheight,keepaspectratio]{../doc/images/animacion-jarra-de-agua-2.png}\\
    \small Jarra, chorro y nivel del agua
  \end{minipage}
  \vspace{0.35cm}

  Las pistas exportadas conservan el movimiento diseñado en Blender y se reproducen dentro del motor OpenGL.
\end{frame}

\begin{frame}{Animaciones por huesos}
  \centering
  \captura{../doc/images/animacion-visitante-pone-ofrenda.png}{Visitante que deposita un pan de muerto}\hfill
  \captura{../doc/images/animacion-visitante-tumba-bailando.png}{Visitante que recorre y baila en el panteón}\\[0.25cm]
  \captura{../doc/images/animacion-esqueleto-bailando-1.png}{Esqueleto: primera pose}\hfill
  \captura{../doc/images/animacion-esqueleto-bailando-2.png}{Esqueleto: segunda pose}
\end{frame}

\begin{frame}{Optimización y compatibilidad}
  \small
  \begin{columns}[T,onlytextwidth]
    \column{0.49\textwidth}
      \begin{block}{Rendimiento}
        Instanciado, descarte por cámara, selección de luces cercanas, caché de materiales y reducción de modelos pesados conservaron el aspecto del altar con menor costo de renderizado.
      \end{block}
    \column{0.48\textwidth}
      \begin{block}{Multiplataforma}
        CMake mantiene un solo proyecto para macOS, Windows y Linux. Rutas portables y stb\_image resolvieron diferencias en carga de recursos.
      \end{block}
      \begin{block}{Distribución}
        GitHub Actions compila Windows x64 y CPack con NSIS produce el instalador junto con recursos, licencia y manual de uso.
      \end{block}
  \end{columns}
\end{frame}

\begin{frame}{Resultados}
  \begin{columns}[T,onlytextwidth]
    \column{0.45\textwidth}
      \includegraphics[width=\linewidth,height=0.65\textheight,keepaspectratio]{../doc/images/final-altar.png}
    \column{0.52\textwidth}
      \begin{itemize}
        \item Escena recorrible con identidad visual coherente.
        \item Phong, Fresnel, skybox e iluminación dinámica.
        \item Nueve animaciones: código, Blender y huesos.
        \item Interacción mediante teclado y ratón.
        \item Instalador autónomo para Windows x64.
        \item Documentación técnica y de usuario en español e inglés.
      \end{itemize}
  \end{columns}
\end{frame}

\begin{frame}{Entrega y demostración}
  \begin{columns}[T,onlytextwidth]
    \column{0.52\textwidth}
      \href{https://youtu.be/2iYnnTLFxNo}{\includegraphics[width=\linewidth,height=0.57\textheight,keepaspectratio]{../doc/images/miniatura.png}}
      \centering\small Video demostrativo: \url{https://youtu.be/2iYnnTLFxNo}
    \column{0.45\textwidth}
      \begin{block}{Repositorio y descarga}
        \href{https://github.com/olveratrech/proyecto-final-cempasuchil}{Repositorio del proyecto}

        \vspace{0.25cm}
        \href{https://github.com/olveratrech/proyecto-final-cempasuchil/releases/latest}{Release más reciente}
      \end{block}
      \small
      \begin{itemize}
        \item Código fuente y modelos de Blender
        \item Reportes finales en ambos idiomas
        \item Instalador y suma SHA-256
      \end{itemize}
  \end{columns}
\end{frame}

\begin{frame}{Conclusiones}
  \begin{itemize}
    \item El reto principal fue integrar modelos, materiales, luces y animaciones creadas con herramientas diferentes dentro de una sola escena estable.
    \item La optimización permitió conservar la iluminación y la densidad visual sin eliminar los elementos que dan identidad a la ofrenda.
    \item Las pruebas multiplataforma mostraron que una entrega completa también requiere rutas portables, empaquetado, documentación y control de versiones.
  \end{itemize}
  \vfill
  \begin{center}
    {\Large\color{cempasuchil}\textbf{Gracias}}\\[0.35cm]
    \small Basilio Illescas Andrés · Cruz Macedo Samuel Santiago\\
    Pérez Olvera Alexis Abraham · Trejo Olvera Emmanuel
  \end{center}
\end{frame}

\end{document}
